\begin{textsample}{POS Dim 3 – vem\_ed\_05 – Score 1.00 – vem\_ed\_05/t019.txt}  \label{ex:f3_pos_040}
Intercâmbios Virtuais: pesquisa 2020 Ao término dos Projetos Colaborativos Internacionais realizados no segundo semestre de 2020, estudantes e professores responsáveis pelos Intercâmbios Virtuais nas 20 Fatecs participantes foram convidados a responder a uma pesquisa de percepção.
Dos 34 professores, 25 responderam; dos 833 alunos, aproximadamente 400 retornaram o questionário.
A seguir, uma síntese dos principais resultados.
ALUNOS 79 \% indicariam a experiência para outros alunos 84 \% consideram que a competência no idioma estrangeiro melhorou 72 \% avaliam interação com colegas estrangeiros como"ótima"ou"boa"86 \% avaliam interação com colegas brasileiros"ótima"ou"boa"63 \% acham que o Intercâmbio virtual melhora desempenho acadêmico PROFESSORES 100 \% recomendariam a experiência para outros docentes 88 \% afirmam que a competência no idioma estrangeiro melhorou. 88 \% avaliam como"ótima"e 12 \% como"boa"a interação com colegas estrangeiros. 92 \% observaram que a interação dos alunos melhorou com a inserção dos projetos colaborativos nas atividades de aula 72 \% concordam totalmente que as colaborações contribuem para pesquisas na área do professor 92 \% concordam totalmente que o intercâmbio virtual possibilita novas atividades acadêmicas a partir da experiência do parceiro internacional.
Ferramentas digitais mais utilizadas Whatsapp \textbf{E-mail} Zoom

% matched lemmas: e-mail
\end{textsample}
